\section{Phép dựng giàn: giấu thiên lệch sau phân phối nhiễu}
\label{sec:ch14-phep-dung-gian}

\index{phép dựng giàn|(}%
Mục trước đã đặt quan hệ thứ nhất vào chỗ của nó: lời giải thích là nạn nhân.
Phép tấn công dựng nên quan hệ ấy có tên là \tn{phép dựng giàn}{scaffolding},
và nó do Slack cùng cộng sự đưa ra~\autocite{mfooling}. Bài ấy không nằm trong
danh mục ba mươi hai bài của sách. Sách vẫn đọc nó trực tiếp và cấp cho nó một
khóa trích dẫn riêng, vì hai mục tiếp theo lập luận trên chính phép dựng ấy, và
bài duy nhất của danh mục có nói tới nó, bài của
mục~\ref{sec:ch14-mot-phong-thu}, là một bài lặp lại thí nghiệm của nó rồi mở
rộng. Đây là lần thứ hai sách trích ra ngoài danh mục, sau bảy khóa mà
chương~\ref{ch:do-do-trung-thuc} phải thêm vì không bài nào trong danh mục định
nghĩa một \tn{chỉ số trung thực}{faithfulness metric}.

\index{phép dựng giàn!tình huống}%
Tình huống mà bài đặt ra là tình huống của người kiểm tra. Một bên nào đó có mô
hình $f$ dùng cho một quyết định quan trọng, chẳng hạn cho tại ngoại hay cho
vay, và mô hình ấy thiên lệch. Khách hàng và cơ quan quản lý chỉ có quyền truy
cập \tn{hộp đen}{black box}: họ đưa đầu vào, nhận đầu ra, rồi chạy LIME hoặc
SHAP để xem mô hình dựa vào đâu. Nếu lời giải thích nêu ra một thuộc tính được
bảo vệ thì mô hình không được duyệt. Mục tiêu của bên triển khai là giữ nguyên
$f$ mà lời giải thích không nêu thuộc tính ấy.

Cái khe hở mà phép tấn công chui qua đã được chương~\ref{ch:lime} và
chương~\ref{ch:shap} mô tả đủ, tuy không ai gọi nó là khe hở. Cả LIME lẫn SHAP
đều không đọc mô hình; chúng chỉ hỏi mô hình. LIME sinh một tập mẫu quanh $x$
bằng cách thêm nhiễu, rồi khớp một mô hình tuyến tính lên các câu trả lời
(mục~\ref{sec:ch03-sinh-lan-can}); SHAP thay các đặc trưng vắng mặt bằng giá
trị lấy từ một
\tn{phân phối nền}{background distribution} rồi hỏi mô hình về những đầu vào lai
ấy (mục~\ref{sec:ch04-gia-tri-lien-minh-va-phan-phoi-nen}). Điểm chung là cả
hai kết luận về hành vi của mô hình trên dữ liệu thật bằng cách hỏi nó về những
điểm không phải dữ liệu thật.

\index{đầu vào ngoài phân phối!làm chỗ dựa cho tấn công}%
Bài kiểm chứng rằng khe hở ấy rộng bằng một thí nghiệm nhỏ, và thí nghiệm ấy là
bước đầu tiên của phép tấn công chứ không phải một minh họa. Các tác giả lấy
tập COMPAS, sinh nhiễu quanh từng điểm theo đúng cách LIME sinh, gộp dữ liệu
thật với dữ liệu nhiễu rồi chiếu cả hai xuống hai thành phần chính đầu tiên.
Hai đám mây tách nhau rõ~\autocite{mfooling}. Nói cách khác, một mẫu do LIME
sinh ra không giống một điểm dữ liệu thật, và phân biệt chúng là một bài toán
phân loại dễ. Chỗ này đã có tên trong sách: đó là \tn{đầu vào ngoài phân
phối}{out-of-distribution input} mà mục~\ref{sec:ch08-huan-luyen-lai} đã gặp
khi đọc phản đối của Hase và cộng sự nhắm vào các họ chỉ số. Ở đó, tính ngoài
phân phối là một nguồn sai lệch mà người chạy chỉ số phải chịu. Ở đây nó là chỗ
dựa: cái gì phân biệt được thì cũng lợi dụng được.

Phần còn lại là một dòng. Bên tấn công giữ mô hình thiên lệch $f$, dựng thêm
một mô hình vô hại $u$ chỉ dựa vào những \tn{đặc trưng}{feature} không tương
quan với thuộc tính được bảo vệ, rồi công bố mô hình

\begin{equation}
\label{eq:ch14-phep-dung-gian}
e(x) =
\begin{cases}
f(x), & \text{nếu } x \in \mathcal{X}_{\mathrm{dist}},\\
u(x), & \text{nếu không,}
\end{cases}
\end{equation}

trong đó $\mathcal{X}_{\mathrm{dist}}$ là tập những điểm dữ liệu thật, tức
những đầu vào mà mô hình sẽ gặp khi chạy~\autocite{mfooling}. Ai quyết định một
đầu vào có nằm trong tập ấy hay không là câu hỏi của đoạn sau. Bài viết mô hình
vô hại là $\psi$; sách viết $u$, vì $\psi_i(x)$ ở phụ lục~\ref{app:ky-hieu} đã
thuộc về \tn{attribution dữ liệu}{data attribution} của
chương~\ref{ch:khung-hop-nhat}.

\index{phép dựng giàn!bộ phát hiện ngoài phân phối}%
Phân biệt hai nhánh trong phương trình~\ref{eq:ch14-phep-dung-gian} là việc của
một bộ phân loại phụ mà bài gọi là
\tn{bộ phát hiện ngoài phân phối}{out-of-distribution detector}. Nó được huấn
luyện theo cách trực tiếp nhất có thể: lấy tập dữ liệu thật, sinh nhiễu quanh
từng điểm theo đúng cách LIME sinh, gán nhãn \enquote{thật} cho tập gốc và
\enquote{nhiễu} cho tập sinh, rồi khớp một mô hình có sẵn lên hai nhãn ấy. Trong
thí nghiệm của bài, mô hình ấy là một rừng ngẫu nhiên một trăm
cây~\autocite{mfooling}. Bộ phát hiện càng chính xác thì phương
trình~\ref{eq:ch14-phep-dung-gian} càng định tuyến đúng, nên chất lượng của nó
là tham số duy nhất quyết định phép tấn công thành công tới đâu, và
mục~\ref{sec:ch14-mot-phong-thu} sẽ đọc đúng đường cong ấy.

Điều đáng nhìn trong hình~\ref{fig:ch14-phep-dung-gian} là hai luồng câu hỏi đi
vào bộ phát hiện ấy rồi ra ở hai nhánh khác nhau, và nhất là mũi tên chấm chấm
từ người kiểm tra tới nhánh thiên lệch: đó là câu hỏi mà không ai đặt được.

\begin{figure}[htbp]
\centering
\begin{tikzpicture}[
  every node/.style={font=\footnotesize, align=center},
  box/.style={draw=black!65, rounded corners=2pt, fill=black!5,
    minimum width=30mm, minimum height=11mm},
  routebox/.style={box, draw=black!85, line width=1pt, fill=black!12,
    minimum width=32mm},
  branchbox/.style={box, fill=black!16, minimum width=26mm},
  freechoice/.style={draw=black!55, dashed, thick, rounded corners=2pt,
    fill=white, minimum width=30mm, minimum height=11mm},
  flow/.style={-{Stealth[length=2mm]}, thick, black!70},
  never/.style={-{Stealth[length=2mm]}, thick, black!45, dotted},
  lbl/.style={font=\scriptsize, align=center, black!55}
]
  \node[box]        (user)    at (0,1.9)    {Người dùng cuối\\hỏi điểm thật $x$};
  \node[box]        (audit)   at (0,-1.0)   {Người kiểm tra\\chạy LIME hoặc SHAP};
  \node[freechoice] (perturb) at (0,-3.3)   {Phân phối nhiễu:\\lựa chọn tự do};

  \node[routebox]   (route)   at (5.4,0.45) {Bộ phát hiện\\ngoài phân phối};

  \node[branchbox]  (f)       at (10.4,1.9) {$f$: thiên lệch};
  \node[branchbox]  (u)       at (10.4,-1.0){$u$: vô hại};

  \draw[flow] (user)    -- (route);
  \draw[flow] (audit)   -- (route);
  \draw[flow] (perturb) -- (audit);
  \draw[flow] (route)   -- (f);
  \draw[flow] (route)   -- (u);

  \node[lbl] at (7.9,1.6)  {điểm thật};
  \node[lbl] at (7.9,-0.8) {mẫu đã nhiễu};

  \draw[never] (audit.west) -- (-2.4,-1.0) -- (-2.4,3.3) -- (10.4,3.3)
    -- (f.north);
  \node[lbl] at (4.0,3.6) {câu hỏi không ai đặt được};
\end{tikzpicture}

\caption{Hai luồng câu hỏi không nhận hai câu trả lời khác nhau; chúng bị
phân biệt rồi định tuyến, và mỗi luồng chỉ gặp một nhánh. Phân phối nhiễu vẽ
nét đứt vì nó là lựa chọn tự do thứ hai trong ba lựa chọn của
mục~\ref{sec:ch08-vong-lap-chung}, và phép tấn công dựng trên đúng lựa chọn
ấy.}
\label{fig:ch14-phep-dung-gian}
\end{figure}

\index{phép dựng giàn!kết quả}%
Kết quả thì rõ tới mức không cần một bảng nào. Bài chạy trên ba tập dữ liệu,
đều là những tập mà thuộc tính được bảo vệ đã biết trước: COMPAS với 6172 bị
cáo, Communities and Crime với 1994 cộng đồng, và German Credit với 1000 người
vay~\autocite{mfooling}. Mô hình thiên lệch được dựng cực đoan có chủ ý, tức là
nó dự đoán chỉ bằng thuộc tính được bảo vệ và không dùng gì khác. Với LIME và
một đặc trưng đánh lạc hướng, lời giải thích chuyển đặc trưng quan trọng nhất
từ thuộc tính được bảo vệ sang đặc trưng đánh lạc hướng ở 100 phần trăm số điểm
kiểm tra của COMPAS và của Communities and Crime. Ở German Credit thì tỉ lệ là
91 phần trăm, và đặc trưng nhận chỗ ấy không phải một đặc trưng dựng thêm mà là
một đặc trưng có sẵn của bộ dữ liệu, tỉ lệ tiền vay trên thu
nhập~\autocite{mfooling}. Với SHAP, cùng thiết lập ấy cho 100 phần trăm ở
Communities and Crime, 84 phần trăm ở COMPAS và 85 phần trăm ở German
Credit~\autocite{mfooling}.

\index{độ chính xác cục bộ!cản phép tấn công}%
Con số đáng đọc kỹ lại là con số nhỏ nhất. Khi bên tấn công dùng hai đặc trưng
đánh lạc hướng thay vì một, SHAP không còn đẩy được thuộc tính được bảo vệ ra
khỏi vị trí thứ nhất một cách đáng tin: tỉ lệ tụt xuống 67 phần trăm với COMPAS
và 71 phần trăm với Communities and Crime~\autocite{mfooling}. Lý do mà bài đưa
ra không phải một chi tiết cài đặt mà là một tiên đề, và là tiên đề mà
chương~\ref{ch:shap} đã đọc kỹ. \tn{Độ chính xác cục bộ}{local accuracy} trong
định nghĩa~\ref{def:ch04-ba-tinh-chat} buộc tổng các điểm số bằng hiệu giữa dự
đoán tại $x$ và dự đoán trung bình dưới phân phối nền, nên khi không có một đặc
trưng nào nổi hẳn lên thì SHAP phải rải điểm cho nhiều đặc trưng thay vì dồn
cho một~\autocite{mfooling}. Việc rải ấy giữ cho thuộc tính được bảo vệ còn sót
lại trong nhóm đầu bảng.

Đây là một tính chất tốt sinh ra như tác dụng phụ. Không ai đặt độ chính xác
cục bộ vào SHAP để chống lại một phép tấn công chưa tồn tại lúc tiên đề được
viết, và bài cũng không nói rằng SHAP an toàn hơn LIME theo nghĩa nào ngoài
nghĩa hẹp này. Điều nó cho thấy là chỗ chịu lực của một phương pháp attribution
có thể nằm ở một ràng buộc toán học chứ không ở một biện pháp phòng thủ, và
mục~\ref{sec:ch14-mot-phong-thu} sẽ đọc một biện pháp phòng thủ được thiết kế
hẳn hoi để so.

\index{phép dựng giàn!độ khớp của mô hình đối kháng}%
Còn một con số nữa mà phép tấn công phải đạt thì mới có nghĩa, và nó là con số
dễ bỏ qua. Mô hình $e$ được công bố phải hành xử đúng như $f$ trên dữ liệu
thật, vì nếu không thì bên tấn công đã đánh mất chính cái thiên lệch mà họ muốn
giữ. Bài đo tỉ lệ hai mô hình cho cùng dự đoán trên tập kiểm tra: với các mô
hình đối kháng nhắm vào LIME thì tỉ lệ ấy là 100 phần trăm trên cả ba tập; với
các mô hình nhắm vào SHAP thì là 91 phần trăm ở COMPAS, 80 phần trăm ở German
Credit và 75 phần trăm ở Communities and Crime~\autocite{mfooling}. Đại lượng
ấy là độ khớp theo nghĩa dòng đầu bảng~\ref{tab:ch01-bon-quan-he}, đo giữa hai
mô hình chứ không giữa một lời giải thích và một mô hình. Bài dùng chính chữ
\enquote{fidelity} cho nó.

\index{phép dựng giàn!bền với siêu tham số}%
Điều cuối cùng của mục này là điều khiến phép tấn công khó gạt đi như một hiện
tượng của một cách cài đặt. Bài thử đổi các \tn{siêu tham số}{hyperparameter}
mà người chạy LIME và SHAP hay đổi, và phép tấn công không suy chuyển. Độ rộng
\tn{hàm nhân}{kernel} của LIME chạy từ $0.15\sqrt{M}$ tới $0.95\sqrt{M}$ quanh
mặc định $0.75\sqrt{M}$, và các phép tấn công vẫn hiệu quả như nhau; đổi chuẩn
khoảng cách từ $\ell_2$ sang $\ell_1$ cũng không đổi gì; với SHAP, đổi số cụm
của phân phối nền từ mười sang năm hoặc mười lăm cũng vậy~\autocite{mfooling}.
Một trường hợp thì đi ngược lại và đáng ghi: dùng phân phối nền toàn số không
làm phép tấn công mạnh lên~\autocite{mfooling}, tức là lựa chọn nền mà
mục~\ref{sec:ch04-gia-tri-lien-minh-va-phan-phoi-nen} đã gọi là chưa được kiểm
định lại là lựa chọn đẩy người kiểm tra vào chỗ dễ bị lừa hơn.

Bảy lựa chọn tự do mà mục~\ref{sec:ch03-cac-tham-so-tu-do} đếm được trong LIME
vì thế không phải bảy chỗ để chỉnh cho phép tấn công hết tác dụng. Chúng nằm
bên trong vùng mà phép tấn công đã bao trọn, bởi phép tấn công không nhắm vào
một giá trị tham số nào mà nhắm vào điều đúng với mọi giá trị: rằng những điểm
được hỏi không phải những điểm mà mô hình sẽ gặp khi chạy.

\index{phép dựng giàn|)}%
